\documentclass[12pt]{article}
\usepackage[T1]{fontenc}
\usepackage[utf8]{inputenc}
\usepackage{lmodern}
\usepackage{textcomp}
\usepackage{enumitem}
\usepackage{geometry}
\geometry{margin=1in}
\begin{document}
\begin{center}
Answer Key - History - K-12 Level Assessment 4 \
\small (Answers are ordered exactly as in the question paper.)
\end{center}

\section*{Multiple Choice}
\begin{enumerate}[label=\arabic*.]
\item \textbf{Answer:} D
\\ \textit{Options:} A) South America.; B) Mesopotamia.; C) the Indus Valley.; D) China.
\\ \textit{Note:} The correct answer is China because archaeological evidence indicates that the earliest known bronze artifacts, dating back to around 3000 BCE, were discovered in ancient Chinese sites such as the Erlitou culture.
\item \textbf{Answer:} A
\\ \textit{Options:} A) Homo erectus.; B) Homo habilis.; C) Homo ergaster.; D) Australopithecus afarensis.
\\ \textit{Note:} Homo floresiensis is considered similar or related to Homo erectus due to their shared characteristics as members of the genus Homo and the evidence suggesting that Homo floresiensis may have evolved from a small-bodied lineage of Homo erectus that adapted to the island environment.
\item \textbf{Answer:} D
\\ \textit{Options:} A) in western Europe, about 3,500 years ago; B) in sub-Saharan Africa, about 8,500 years ago; C) in North America, about 9,500 years ago; D) in the Middle East, about 10,500 years ago
\\ \textit{Note:} Archaeological evidence, including the analysis of ancient sites and remains, suggests that the domestication of cattle began in the Middle East approximately 10,500 years ago, marking a significant development in early agricultural societies.
\item \textbf{Answer:} A
\\ \textit{Options:} A) monumental works; B) burials with grave goods; C) leisure and a sedentary lifestyle; D) all of the above
\\ \textit{Note:} Monumental works, such as large buildings, monuments, or infrastructures, indicate a complex society because they require advanced organization, resources, and labor, reflecting social structure, cultural values, and technological capabilities.
\item \textbf{Answer:} A
\\ \textit{Options:} A) Anthropology is holistic and integrative in its approach.; B) Anthropology is less holistic and more theoretical in its approach.; C) Anthropology tends to focus on specific topics in history.; D) Anthropology is more narrowly focused than the other social sciences.
\\ \textit{Note:} The correct answer is based on the key concept that anthropology studies human behavior and culture through a comprehensive lens, integrating insights from various disciplines, while other social sciences like economics and sociology often focus on specific aspects of social life in a more segmented manner.
\end{enumerate}

\section*{True / False}
\begin{enumerate}[label=\arabic*.]
\setcounter{enumi}{5}
\item \textbf{Answer:} False
\\ \textit{Options:} A) True; B) False
\\ \textit{Note:} The world's first civilization actually developed in Mesopotamia, located between the Tigris and Euphrates rivers, rather than in Eastern China.
\item \textbf{Answer:} True
\\ \textit{Options:} A) True; B) False
\\ \textit{Note:} Brachiation refers specifically to the locomotion method used by certain primates, such as gibbons and spider monkeys, which involves swinging from branch to branch using their arms, confirming the statement as true.
\item \textbf{Answer:} True
\\ \textit{Options:} A) True; B) False
\\ \textit{Note:} Homo erectus is considered a stable and long-lived species because it existed for nearly 1.9 million years, exhibiting adaptability and a wide geographic range that allowed it to thrive during a critical period in human evolution.
\item \textbf{Answer:} False
\\ \textit{Options:} A) True; B) False
\\ \textit{Note:} Cahokia was known for having over 100 burial mounds and a population that may have reached up to 20,000 at its peak, making the statement false.
\item \textbf{Answer:} False
\\ \textit{Options:} A) True; B) False
\\ \textit{Note:} The Mogollon culture was primarily located in the southwestern United States, particularly in present-day New Mexico and Arizona, not in the Ohio River valley.
\end{enumerate}

\section*{Long Form Response}
\begin{enumerate}[label=\arabic*.]
\setcounter{enumi}{10}
\item \textbf{Answer:} The first evidence of metal use in South America dates back to around 3100 years before the present (B.P.). This significant development marked a crucial turning point in the region's history, as it allowed for advancements in tools and weapons, which in turn facilitated agricultural practices and trade. The introduction of metalworking not only improved the efficiency of everyday tasks but also influenced the social structures within communities, leading to increased specialization and the emergence of more complex societies. As cultures began to create intricate metal artwork and tools, this contributed to a rich cultural heritage that shaped the identity of various South American civilizations. Overall, the early use of metal played a vital role in the technological and cultural evolution of the region.
\\ \textit{Note:} This answer is correct because it identifies the date of the first metal use in South America as a pivotal moment that spurred technological advancements, improved agricultural practices, enhanced trade, and contributed to social stratification, all of which significantly impacted the region's development and culture.
\item \textbf{Answer:} Hominids evolved to walk on two legs for several reasons, including the need to travel long distances for food, the ability to use tools, and the potential for braincase expansion. Bipedalism freed the hands for tool use and carrying objects, while also allowing hominids to see over tall grass and spot predators. One hypothesis that is least supported by evidence is the idea that bipedalism primarily evolved to regulate body temperature, as this does not fully account for the anatomical changes observed in early hominids. Overall, the expansion of the braincase is a significant factor as it correlates with the development of higher cognitive functions, which were crucial for survival and adaptation.
\\ \textit{Note:} correct because it identifies key reasons for the evolution of bipedalism in hominids, such as tool use and predator awareness, while also accurately highlighting the body temperature regulation hypothesis as the least supported due to insufficient anatomical evidence.
\end{enumerate}

\vfill
\noindent\textit{End of Answer Key}
\end{document}